% notation_table.tex -- the ONE notation table shared by Paper W and the companion (version 3.0).  Edit here only.
% [W] = Paper W (Tilings of W_M and T_M below the threshold scale); [C] = the companion paper.
\begin{center}\small
\begin{tabular}{@{}p{2.6cm}p{9.9cm}@{}}
\toprule
symbol & meaning ([W] = Paper W; [C] = the companion paper)\\
\midrule
\multicolumn{2}{@{}l}{\emph{used in both papers}}\\
$e<f$ & coprime positive integers defining the tile; $d=f-e$\\
$a,b,c,\kappa$ & $a=ef$, $b=f^2-e^2$, $c=f^2$, $\kappa=b+c=2f^2-e^2$: the sides of the tile $\tile(e,f)$ and the sum $b+c$\\
$\alpha,\beta,\gamma$ & angles of the tile opposite $a,b,c$; $3\alpha+2\beta=\pi$, $\gamma=2\alpha+\beta$\\
scale $s$ & a triangle with sides $sa,sb,sc$ has \emph{scale} $s$ and is tiled by $s^2$ tiles (ladder subdivision, [W])\\
$W_M$, $T_M$ & the two targets of [W]; $M$ is the \emph{scale} of the target; $W_M$ is tiled by $M^2(2f^2-e^2)$ tiles, $T_M$ by $M^2(3f^2-e^2)$\\
$M^*(e,f)$ & threshold scale $\max(\lceil f/e\rceil,\lceil f/(f-e)\rceil+1)$ of Theorem~W; (H1), (H2) are $Me\ge f$ and $(M-1)(f-e)\ge f$, together equivalent to $M\ge M^*(e,f)$\\
\midrule
\multicolumn{2}{@{}l}{\emph{used in [W] only}}\\
$n,\rho,\delta$ & $n=\lfloor eM/f\rfloor$, $\rho=eM-fn$, $\delta=1$ if $\rho=0$ and $0$ otherwise\\
$t,r,q$ & the three free integers of the six-piece layout; canonical choice $t=n-\delta$, $r=n+1$, $q=n$\\
$\pi_0,\dots,\pi_5$ & the six pieces of the layout; $O,A,B,P,Q,R,\dots$ its points\\
\midrule
\multicolumn{2}{@{}l}{\emph{used in [C] only}}\\
$B_g$ & the isosceles target with base angles $\alpha$; $g$ is its scale; $b\kappa g^2$ tiles\\
row $(e,f)$ & the tile $\tile(e,f)$; $G(e,f)$ is the set of scales $g$ for which $B_g(e,f)$ is tiled\\
$F_0,\dots,F_8$ & the families of Conjecture~C\\
\bottomrule
\end{tabular}
\end{center}
